\documentclass[11pt]{article}
\usepackage{amsmath,amssymb,amsthm}
\usepackage{booktabs}
\usepackage[margin=1in]{geometry}
\newtheorem{theorem}{Theorem}
\newtheorem{lemma}[theorem]{Lemma}
\title{Problem 782: Distinct Rows and Columns}
\author{Project Euler Solutions}
\date{}
\begin{document}
\maketitle

\section{Problem Statement}
A binary matrix is \emph{distinct} if all rows are different and all columns are different. Let $D(n,m)$ be the number of distinct $n \times m$ binary matrices. Find $\sum_{m=1}^{20} D(20, m)$ modulo $1\,000\,000\,007$.

\section{Mathematical Foundation}

\begin{theorem}[Row Distinctness]
The number of $n \times m$ binary matrices with all $n$ rows distinct is the falling factorial
\[
R(n,m) = (2^m)_n = 2^m(2^m - 1)(2^m - 2)\cdots(2^m - n + 1)
\]
when $2^m \ge n$, and $0$ otherwise.
\end{theorem}

\begin{proof}
Each row is a binary vector in $\{0,1\}^m$, giving $2^m$ possible row values. Choosing $n$ distinct rows in order is the falling factorial $(2^m)_n$.
\end{proof}

\begin{lemma}[Column Partition Reduction]
The computation reduces to:
\[
D(n,m) = \sum_{p=0}^{m} (-1)^{m-p}\, S(m,p)\, (2^p)_n
\]
where $S(m,p)$ is the Stirling number of the second kind (partitioning $m$ columns into $p$ nonempty groups), and $(2^p)_n$ is the falling factorial.
\end{lemma}

\begin{proof}
Apply M\"obius inversion on the partition lattice of columns. When columns are partitioned into $p$ equivalence classes (identical columns), the effective matrix has $p$ distinct column-types. The row-distinctness count for binary vectors of length $p$ is $(2^p)_n$. Summing over all partitions with the appropriate sign from inclusion-exclusion over column coincidences yields the stated formula.
\end{proof}

\begin{theorem}[Full Formula]
Let $D(n,m)$ count $n \times m$ binary matrices where all rows are distinct and all columns are distinct. Then
\[
D(n,m) = \sum_{p=0}^{m} (-1)^{m-p}\, S(m,p)\, (2^p)_n.
\]
\end{theorem}

\begin{proof}
Combine Theorem~1 and Lemma~1. The inclusion-exclusion over column equalities partitions the column set. Each partition into $p$ parts contributes $S(m,p)$ terms with sign $(-1)^{m-p}$ (from the M\"obius function of the partition lattice), and each such configuration admits $(2^p)_n$ row-distinct matrices.
\end{proof}

\section{Algorithm}

\begin{verbatim}
function D(n, m):
    Precompute Stirling numbers S(m, p) for p = 0..m
    using S(i,p) = p*S(i-1,p) + S(i-1,p-1)

    result = 0
    for p = 0 to m:
        ff = falling_factorial(2^p, n) mod MOD
        sign = (-1)^(m - p)
        result += sign * S[m][p] * ff
    return result mod MOD

function solve():
    total = sum of D(20, m) for m = 1 to 20
    return total mod 10^9+7
\end{verbatim}

\section{Complexity Analysis}

\begin{itemize}
    \item \textbf{Time:} $O(m^2)$ per call to $D(n,m)$ for Stirling numbers, plus $O(m \cdot n)$ for falling factorials. Total over all $m$: $O(m_{\max}^3 + m_{\max}^2 \cdot n)$.
    \item \textbf{Space:} $O(m^2)$ for the Stirling number table.
\end{itemize}

\section{Answer}

\[
\boxed{318313204}
\]

\end{document}
